\documentclass[a4paper]{article}
\usepackage[warn]{mathtext}
\usepackage{graphicx} % Required for inserting images
\usepackage{wrapfig}
\usepackage{enumitem}
\usepackage[center]{caption}
\usepackage[english, russian]{babel}
\usepackage{amsmath}
\usepackage{geometry}
\setcounter{totalnumber}{5}

\geometry{verbose, a4paper, tmargin=2cm, bmargin=2cm, lmargin=2.0cm, rmargin=2.0cm}

\begin{document}

\begin{titlepage}
	\centering
	{\scshape\Large МОСКОВСКИЙ ФИЗИКО-ТЕХНИЧЕСКИЙ ИНСТИТУТ 
	(НАЦИОНАЛЬНЫЙ ИССЛЕДОВАТЕЛЬСКИЙ УНИВЕРСИТЕТ) % название ВУЗа большим шрифтом
    Физтех-школа аэрокосмических технологий}
	
	\vspace{4cm} % отступ 4 см по вертикали
	{\LARGE Вопрос по выбору. Механика.}
	\vspace{1cm} % ещё 1 см
	
{\huge\bf Изучение нелинейных колебаний  длиннопериодного маятника}
	
	\vspace{1cm} % ещё 1 см
	\vfill % заполнение по вертикали пустотой (LaTeX сам решит, сколько здесь отступить)
	
    \begin{flushright} % Этот текст будет с правого края страницы
	   {\LARGE Губарев Никита, Б03-502}
    \end{flushright}
    \vfill 
	\today % Дата сборки документа
\end{titlepage}
\tableofcontents 
 \section{Аннотация}
Цель работы: Исследование нелинейных колебаний на примере длиннопериодного маятника. Определение зависимости периода колебаний маятника от амплитуды.

\section{Теоретические сведения}
\subsection{Безразмерное уравнение движения физического маятника с вязким трением}
Уравнение движения физического маятника с учётом вязкого трения:

\begin{center}$I\ddot{\varphi}+b\dot{\varphi}+sin(\varphi)mga =0$ (1)\end{center}
где $I$ - момент инерции, $b$ -коэффициент момента сил вязкого тре-
ния, $m$ - масса маятника, $g$ -ускорение свободного падения, $a$ - рассто-
янии от точки подвеса до центра масс. Введя частоту $\omega_0$, перепишем уравнение (1) в виде:

\begin{center}$\omega_0=\sqrt{\frac{mga}{I}}$\end{center}

\begin{center}$\ddot{\varphi}+\frac{b}{I}\dot{\varphi}+\omega_0^2 sin(\varphi) =0$ (2)\end{center}
И приведем полученное уравнение (2) к безразмерному виду, удобному для анализа:
\begin{center}$\tau=\omega_0t$, $\beta=\frac{b}{\omega_0I}$\end{center}

\begin{center}$\varphi''+\beta\varphi'+sin(\varphi)=0$ (3)\end{center}
В уравнении (3) производные берутся по безразмерному времени $\tau = \omega_0t$. При этом максимальное значение амплитуды безразмерной скорости $\varphi'$ соответствующей колебательному движению маятника без затухания составляет:  $\varphi_{max}'=2$. Отличие решений уравнения (3) от решения нелинейного маятника без затухания определяется только значением безразмерного параметра $\beta$, который в дальнейшем предполагается малым:
\begin{center}$0<\beta<0.1$ (4)\end{center}
\subsection{Фазовая диаграмма нелинейного маятника с затуханием}
С помощью численного решения уравнения (3) можно определить фазовый портрет нелинейного маятника для заданных начальных условий и
параметра затухания $\beta$. В приложении 1 приведён пример фазового портрет
для $\beta=0.1$ и с начальной безразмерной скоростью $\varphi'=2$ . В приложении
также показаны красным цветом сепаратрисы нелинейного маятника без
затухания ограничивающие область колебательных решений.

\subsection{Закон сохранения энергии для нелинейного маятника с затуханием}
Интегрирование уравнения (3) по углу отклонения:
\begin{center}$\frac{\varphi'^2}{2}-cos(\varphi)=-\beta\int\varphi'd\varphi + Const$ (5)\end{center}
Постоянную интегрирования найдём из начальных условий. В начальный момент времени ($\tau=0$) будем считать, что маятник проходит
положение равновесия ($\varphi=0$) , и имеет скорость $\varphi_0'$, тогда:
\begin{center}$Const = \frac{\varphi_0'^2}{2}-1$ (6)\end{center}
В момент первого максимального отклонения маятника:
\begin{center}$-cos(\varphi_x) = -\beta\int\limits_{0}^{\varphi_x}\varphi'd\varphi + \frac{\varphi_0'^2}{2} -1$\end{center}

\begin{center}$\beta\int\limits_{0}^{\varphi_x}\varphi'd\varphi = \frac{\varphi_0'^2}{2} - 2sin^2(\frac{\varphi_x}{2})$ (7)\end{center}
При возвращении к положению равновесия угловая скорость маятника $\varphi_1'$ из-за трения станет меньше чем начальная $\varphi_0'$:

\begin{center}$\varphi_1'^2-1=\frac{\varphi_0'^2}{2}-1-\beta\int\limits_{0}^{\varphi_x}\varphi'd\varphi-\beta\int\limits_{\varphi_x}^{0}\varphi'd\varphi$\end{center}

\begin{center}$\varphi_1'^2=\frac{\varphi_0'^2}{2}-\beta\int\limits_{0}^{\varphi_x}\varphi'd\varphi-\beta\int\limits_{\varphi_x}^{0}\varphi'd\varphi$ (8)\end{center}
В выражения (7) и (8) входят интегралы учитывающие потери энер-
гии на трение. Потери механической энергии связаны с работой сил тре-
ния.

\subsection{Связь апмлитуды колебаний и угловых скоростей в положении равновесия.}
Для оценки значения интегралов в уравнении (8) используем следующее
приближение которое, обосновывается в приложении 2. Интегралы, входящие в (8) равны площадям заштрихованных в приложении 2 областей. Примем следующее приблизительное правило:
\begin{center}$\int\limits_{0}^{\varphi_x}\varphi'd\varphi \approx A\varphi_x\varphi_0'$ \space{} $\int\limits_{\varphi_x}^{0}\varphi'd\varphi \approx -A\varphi_x\varphi_1'$ (9)\end{center}

которое означает, что отношение площади верхней заштрихованной области к площади прямоугольника со сторонами $\varphi_0'$ и $\varphi$ такое же как и отношение нижней заштрихованной области к площади прямоугольника со сторонами $\varphi_1'$ и $\varphi_x$. Знак ‘–‘ во втором уравнении введён для учёта отрицательного значения скорости $\varphi_1'$ (заметим, что оба интеграла положительны). Это приблизительное равенство тем лучше будет выполняться, чем меньше параметр безразмерного затухания. Тогда уравнение (8) преабразуется:
\begin{center}$\frac{\varphi_1'^2}{2}=\frac{\varphi_0'^2}{2}-\beta A\varphi_x(\varphi_0'-\varphi_1') \Rightarrow \varphi_0'+\varphi_1'=2\beta A\varphi_x$ (10)\end{center}
А уравнение (7) приводится к виду:
\begin{center}$\beta A\varphi_x\varphi_0'=\frac{\varphi_0'^2}{2}-2sin^2(\frac{\varphi_x}{2}) \Rightarrow 2\beta A\varphi_x\varphi_0'=\varphi_0'^2-4sin^2(\frac{\varphi_x}{2}) \Rightarrow \varphi_0'^2+\varphi_1'\varphi_0' = \varphi_0'^2-4sin^2(\frac{\varphi_x}{2}) \Rightarrow \varphi_1'\varphi_0'=-4sin^2(\frac{\varphi_x}{2})$\end{center}

\begin{center}$sin^2(\frac{\varphi_x}{2})= -\frac{\varphi_1'\varphi_0'}{4}$ (11а)\end{center}
Здесь знак ‘-’ в правой части учитывает то, что скорости $\varphi_0'$ и $\varphi_1'$ всегда
должные иметь разные знаки.
Другой способ оценки интегралов (9) состоит в том, что угловая скорость $\varphi'$ в подынтегральном выражении выводится через формулы (5)
и (6) но с параметром затухания $\beta=0$. (Метод последовательных приближений.) Как показано в приложении 3, это приводит к следующему
соотношению между амплитудой и скоростями:
\begin{center}$sin(\frac{\varphi_x}{2})^2= \frac{\varphi_1'^2+\varphi_0'^2}{8}$ (11б)\end{center}
Среднее между (11a) и (11б):
\begin{center}$sin(\frac{\varphi_x}{2})^2= \frac{(\varphi_1' - \varphi_0')^2}{16}$ (11в)\end{center}
Фактически, (11б) – это средний квадрат синуса половинного угла 
амплитуды, (11в) – квадрат среднего синуса половинного угла амплитуды, а (11a) – это среднее геометрическое. (Имеется в виду усреднение
по двум значениям $sin(\frac{\varphi_x}{2})^2$ получаемым из модели без учёта трения,
$\beta=0$, для скоростей $\varphi_0'$ и $\varphi_1'$.) Отметим, что эти выражения позволяют
вычислять амплитуду колебания без явного знания значения параметра затухания. Но информация об этом параметре связана с различием
угловых скоростей $\varphi_0'$ и $\varphi_1'$.

В приложении 4 показаны графики относительной ошибки выражений
(11a), (11б) и (11в) для $\beta=0.1$ в зависимости от амплитуды, полученные на основе численного решения. Видно, что выражение (11в) является наиболее точным для оценки амплитуды колебаний, относительная погрешность не превышает величины 0.2\%. Для значений $\beta<0.1$ формулы
(11в) точность будет ещё выше. 
\subsection{Период колебаний.}
Под половиной периода колебаний маятника с затуханием понимается
время между двумя последовательными моментами времени прохождения маятником положения своего равновесия. Полупериод:
\begin{center}$\tau_{\frac{1}{2}}=\int\limits_{0}^{\varphi_x}\frac{d\varphi}{\varphi'} + \int\limits_{\varphi_x}^{0}\frac{d\varphi}{\varphi'} $ (12)\end{center}
Если полностью пренебречь затуханием, тогда из (5), (6) и (7):
\begin{center}$\int\limits_{0}^{\varphi_x}\frac{d\varphi}{\varphi'}=\int\limits_{0}^{\varphi_x}\frac{d\varphi}{2\sqrt{sin(\frac{\varphi_x}{2})^2-sin(\frac{\varphi}{2})^2}} = K(sin(\frac{\varphi_x}{2})^2) \approx  \frac{\pi}{2}(1+(\frac{sin(\frac{\varphi_x}{2})}{2})^2+\frac{9}{4}(\frac{sin(\frac{\varphi_x}{2})}{2})^4+\frac{25}{4}(\frac{sin(\frac{\varphi_x}{2})}{2})^6+...)\approx  \frac{\pi}{2}(1+(\frac{\varphi_x}{4})^2+\frac{11}{12}(\frac{\varphi_x}{4})^4+...)$ (13)\end{center}
Через K() обозначена специальная математическая функция - полный эллиптический интеграл первого рода. При переходе от безразмерных величин обратно к размерным получим следующее выражение для периода нелинейных колебаний:
\begin{center}$T=2\frac{\tau_{\frac{1}{2}}}{\omega_0}=T_0\frac{2 sin(\frac{\varphi_x}{2})^2}{\pi}\approx  T_0(1+(\frac{\varphi_x}{4})^2+\frac{11}{12}(\frac{\varphi_x}{4})^4+...)$ (14)\end{center}
Здесь $T_0$ – период линейных колебаний маятника. В приложении 5
показан способ численного вычисления полного эллиптического интеграла $K(x^2)$, который используется для теоретического определения значения периода при заданной амплитуде колебаний.

\subsection{Определение периода линейных колебаний.}

Измерение периода линейных колебаний маятника можно проводить двумя разными способами. В первом случае проводится прямое измерение
периода малых колебаний маятника, когда влияние нелинейности пренебрежимо мало. Во втором случае проводится измерение угловой скорости
маятника при прохождении через положение равновесия, при движении
из начального состояния с углом отклонения равным 180 градусов и нулевой начальной скоростью. Период определяется косвенно – через измеренную скорость. Недостаток первого метода состоит в том, что энергия малых колебаний очень мала и поэтому оказывается существенным
влияние сил трения разной природы (сухое трение, флуктуации вязкого
трения о воздух связанные с потоками воздуха через спицы маятника).
Во втором случае энергия маятника значительно больше, хотя потери на
трения всё равно не равны нулю. Рассмотрим второй метод более подробно. При отсутствии трения закон сохранения энергии связывает угловую
скорость $\omega_x$ и максимальную амплитуду качаний маятника, которая равна 180 градусам:
\begin{center}$\frac{I\omega_x^2}{2}=2mga$\end{center}
Обозначения такие же как и в уравнении (1). Выразим отсюда угловую
скорость $\omega_x$ через угловую частоту линейных колебаний маятника $\omega_0$:
\begin{center}$\omega_x=2\sqrt{\frac{mga}{I}}=2\omega_0$ (15)\end{center}
Таким образом, измеряемая угловая скорость получается в два раза
больше чем угловая частота линейных колебаний. Период линейных колебаний:
\begin{center}$T_0=\frac{4\pi}{\omega_x}$ (16)\end{center}
Уравнение (16) справедливо при отсутствии трения, в случае вязкого
трения можно оценить его вклад в качестве поправки. В приложении
6 выведено выражение (25) в котором учтено влияние трения в рамках
первого приближения. Именно это выражение используется в работе при
обработке предварительных измерений.

\subsection{Измерения скорости и полупериода колебаний с помощью оптоэлектронного датчика.}
Оптоэлектронный датчик располагается так, что спица маятника пересекает луч оптопары тогда, когда маятник проходит через положение
равновесия маятника. Спица маятника имеет в месте прохождения луча угловую ширину $\Delta\varphi$; моменты пересечения спицей луча происходят
так, как показано на рисунке. Здесь отмечены моменты нарушения $T1_i$
и восстановления $T2_i$ оптического тракта оптопары ( i = 1, 2, 3..).
Зная угловой размер спицы маятника и учитывая, что этот размер
достаточно мал ($\Delta\varphi \approx 0.005$), можно определять угловую скорость
маятника в моменты прохождения положения равновесия:
\begin{center}$\dot{\varphi_i}=\frac{\Delta \varphi}{\Delta t_i}$, \space $\Delta t_i=T2_i-T1_i$ (17)\end{center}
Полупериод колебания (приложение 7):
\begin{center}$T_{\frac{1}{2}, i} = \frac{T2_i+T1_i}{2}-\frac{T2_{i-1}+T1_{i-1}}{2}$ (18)\end{center}
Эти выражения используются для определения измеряемых в опытах
величин: угловой скорости и полупериода колебаний. Период колебаний
определяется как сумма двух последовательных полупериодов. Соответствующая этому периоду амплитуда – как среднее арифметическое от амплитуд (формула (11в) ) двух последовательных качаний маятника. Если вместо углового размера $\Delta\varphi$ в формуле (17) использовать эффективный угловой размер $\Delta\varphi_n$:
\begin{center}$\Delta\varphi_n=\frac{\Delta\varphi}{\omega_0}$ (19)\end{center}
то угловая скорость будет определяться в безразмерных единицах. В опытах с качанием маятника с начальной амплитудой в 180 градусов угловая скорость в положении равновесия согласно уравнению (15): $\omega_x=2\omega_0$, с другой стороны из (17) и (18) получаем следующее выражение для эффективного углового размера спицы $\Delta\varphi_n$:
\begin{center}$\omega_x=\frac{\Delta\varphi}{\Delta t}=2\omega_0 \Rightarrow 2\Delta t =\frac{\Delta \varphi}{\omega_0}=\Delta \varphi_n$ (20)\end{center}
Формула (20) выполняется точно, если трение отсутствует, при наличии
слабого вязкого трения можно сделать поправку первого приближения.
В приложении 6 показан вывод с учётом этой поправки – выражение
(26).
\section{Методика измерений}
\begin{enumerate}
\item Установить грузы на спицах маятника так, чтобы на верхней спице груз располагался максимально близко к оси качания, а на нижней - на максимальном удалении от оси. Измерить расстояния от грузов до оси качания.
\item Подобрать корректное положение оптопары. 
\item Провести серию предварительных измерений для определения периода линейных колебаний и эффективного углового размера спицы.
\item Провести серию основных измерений для определения зависимости периода от амплитуды колебаний.
\item Повторить пункты 1-4 для других положений нижнего груза.
\item Построить на одном графике зависимости нормализованного периода $T/T_0$ от амплитуды, полученные для различных положений груза. Определить, как при этом меняется характер линейности.
\item Построить зависимость периода $T_0$ линейных колебаний от положения груза. Определить из этого графика массу перемещаемого груза.
\item Проанализировать графики зависимости параметра затухания от средней скорости маятника. Какая модель трения больше подходит в том или ином диапазоне средних скоростей?
\item Построить зависимость усредненного коэффициента от положения груза. Усреднение необходимо делать по области амплитуд колебаний, где коэффициент трения практически не зависит от амплитуды. Объяснить полученный результат
\end{enumerate}
\section{Результаты измерений и обработка данных}

\section{Выводы}


\section{Приложения}
Приложение 1: Пример фазового портрета для нелинейного маятника с зату-
ханием вычисленный с помощью численного интегрирования уравнения
движения (3)

Приложение 2: Участок траектории маятника на фазовой плоскости

Приложение 3: Метод последовательных приближений и вывод гениальной формулы (11б)

Приложение 4: Относительная точность формул (11a), (11б) и (11в)

Приложение 5: способ численного вычисления полного эллиптического интеграла $K(x^2)$

Приложениe 6:  выражение (25) в котором учтено влияние трения в рамках первого приближения

Приложение 7: Зависимость угла поворота маятника от времени – синяя кривая.
Красные кривые - положение граней спицы маятника. Также показаны
моменты времени, регистрируемые с помощью оптопары
\end{document}

